\section{Factoring intermediate components instead of correlating all spins}
\label{sec:factorization}

\subsection{Why the global orbit table can still be wasteful}

A poor crossing order can create several disjoint components of processed
Tait edges. Before any processed edge joins them, their partial weights
factor. A global $S_q$ orbit table nevertheless records whether an active
spin in one component happens to have the same color as a spin in another.
Those relative color identities matter only when a future edge first
connects the two components. Recording them earlier can create many keys
with an exactly factorizable coefficient.

This is an exact independence statement about the processed interaction
graph. It does not follow merely because two groups of vertices appear far
apart in a drawing, nor does it assume a graph separator whose crossing
interactions have been discarded. We keep one tensor for each connected
component of the graph formed by introduced vertices and processed edges.
Connectivity through a forgotten vertex remains connectivity: forgetting
does not undo an interaction already summed into the tensor.

An isolated newly introduced vertex has the one-key table
$\{(0)\mapsto q\}$. An edge internal to a component uses the usual equality
weight and forgetting operation. When no active vertices remain in a
component, its scalar partition factor is multiplied into the accumulated
scalar and its tensor is released. Only an edge joining two distinct
processed components requires a new operation.

\subsection{The relative-color merge}

Let a left pattern have $k$ classes and a right pattern have $\ell$ classes.
Within each pattern, distinct classes must receive distinct colors.
Across the two patterns, an arbitrary partial bijection may identify
some left and right classes. If $j$ pairs are identified, the merged pattern
has $K=k+\ell-j$ classes, and it is possible exactly when $K\le q$.
The number of such alignments is
\begin{equation}
 A_q(k,\ell)=
 \sum_{\substack{0\le j\le\min(k,\ell)\\k+\ell-j\le q}}
 \binom{k}{j}\binom{\ell}{j}j!.
 \label{eq:alignments}
\end{equation}
It is at most $q!$ when $k,\ell\le q$. One way to see this bound is to fix
distinct actual colors on all left classes and inject the alignment choices
into the injective assignments of colors to the $\ell$ right classes by
assigning fresh colors in a prescribed order. There are
$\fall{q}{\ell}\le q!$ such assignments. The implementation enumerates
alignments directly, giving fresh classes their next canonical label.

\begin{theorem}[Aggregate component merge]\label{thm:merge}
Work over $R_q$, or over a field of characteristic zero or greater than $q$.
Suppose the left and right orbit coefficients are $C_L$ and $C_R$.
For one relative alignment with $K$ total classes, its aggregate coefficient
before the joining edge is applied is
\begin{equation}
 \frac{C_L}{\fall{q}{k}}\,
 \frac{C_R}{\fall{q}{\ell}}\,
 \fall{q}{K}.
 \label{eq:merge}
\end{equation}
After multiplication by the joining-edge weight, ordinary forgetting
and canonical aggregation produce the correct merged tensor.
\end{theorem}
\begin{proof}
The orbit of a pattern with $k$ nonempty classes has
$\fall{q}{k}$ labeled active assignments. Color relabeling preserves the
sum over forgotten spins, so all those active assignments have the same
partial value. That value is $C_L/\fall{q}{k}$. The same statement holds
on the right.

Before the joining edge, the two processed components share no vertex or
processed interaction, so their per-assignment partial values multiply.
The aligned combined pattern has $\fall{q}{K}$ labeled assignments.
Multiplying by this orbit size gives \eqref{eq:merge}.
Every combined pattern determines exactly its restrictions and its
partial-bijection alignment, so no term is omitted or counted twice.
The joining edge and restriction to surviving vertices then obey
Proposition~\ref{prop:orbit}.
\end{proof}

In the modular backend, $q<p$ makes every falling factorial a unit in
$\F_p$, and inverses can be precomputed. In the exact backend, each
coefficient is an integral multiple of its orbit size in $R_q$.
Both integer coordinates are therefore divisible by that positive integer.
The implementation checks the remainders and performs exact integer
division. It then multiplies the resulting pairs.
This operation uses no fraction field and no approximate division.
Unlike the one-tensor transfer, it is not literally division-free; the
divisions are certified divisions by known small integers at fixed $q$.

\subsection{A stronger parameter: the largest processed component}

For each step, consider each still-active processed component and count its
active vertices. Define $g$ to be the largest such count over all components
and all \emph{post-edge} steps. Thus $g\le\max_i f_i\le w/2$.

\begin{theorem}[Component-frontier complexity]\label{thm:componentwidth}
At fixed $q$, the factored transfer uses
\[
 n\,\poly(g,\log n)\,q!\,q^{g+2}
\]
coefficient operations and key-management work, with
$O(nq^g)$ logical table keys. Over $R_q$, coefficient bit lengths remain
$O(n\log q)$; the bit-operation bound is obtained by multiplying the
displayed work by a polynomial in $n\log q$.
\end{theorem}
\begin{proof}
An edge can cause only its two incident vertices to be forgotten: any
vertex whose last incident edge is the current edge must be an endpoint.
Therefore a component's pre-forget active set at this step has at most
$g+2$ vertices. For a merge, this set is the union of the two input active
sets, including newly introduced endpoints. If their sizes are $a,b$, then
$a+b\le g+2$. The number of pairs of input keys is at most
$q^a q^b\le q^{g+2}$; each pair has at most $q!$ alignments.
An internal edge has the same bound without the alignment factor.
There are $n$ edges and at most $O(n)$ active components.

For the exact coefficient bound, each component coefficient is a sum over
assignments on a subset of vertices and a product over its processed edges.
The argument of Theorem~\ref{thm:bits} applies componentwise. During a merge,
the division removes an orbit multiplicity, and the final insertion is a
sum over a subset of assignments on the union. Its length is again bounded
by the same global $q^v(q-1)^n$ envelope.
\end{proof}

Bounding each input component separately by $g$ and then multiplying two
$q^g$ tables would give the weaker exponent $2g$.
The fact that only the current edge's endpoints can be forgotten is what
gives the stronger $g+2$ bound. A measurement of the largest table size alone
would not prove either complexity result: merge Cartesian products and
alignment work must be counted.

\subsection{Implementation and remaining limitations}

The package supplies modular and exact factored kernels. Their state budget
bounds the sum of represented keys over current component tables. During an
update, the old input tables and their replacement may coexist; normalized
merge lists and cached alignment patterns also occupy memory.
Consequently that counter is a logical state budget, not an allocation or
resident-memory ceiling. Transition counts include each relative alignment
that is actually processed.

At $q=5$ there are at most $120$ alignments per input-key pair; at $q=6$
there are at most $720$. These constants can make factorization slower when
components repeatedly merge without having saved many keys. The measured
negative controls in Section~\ref{sec:experiments} are therefore essential.
The implementation does not factor a connected tensor merely because its
coefficient array happens to have low rank, and it does not split components
again when active graph connectivity changes. Such extensions require
additional algebra and are research topics rather than current guarantees.
